\section{The Model}
\label{sec:model}

\subsection{Cells, paths and mode crossing}
\label{sec:model-cells}

A beni value at run time is a JavaScript value. A \emph{cell} $c$ is one JavaScript array that
backs a \code{List} value; two list values that are the same cell are \code{===}, and editing a cell
in place changes what every reference to it sees. A list value is a cell or a \emph{view} of one, a
header over a base array. A view is a holder of its base, and writing through a view writes the base
(\S\ref{sec:hard-views}).

A \emph{path} is a chain of steps from a value into it: a record field $.f$, a tuple index $.i$, a
constructor argument $C\#i$, or a list position $[\ast]$, which stands for every element. Paths are
cut at $k$ steps ($k = 8$ in our design); a value deeper than the cut is approximated by $\top$
below. A path is \emph{list-typed} when the solved type at its end is \code{List a}, and only
list-typed paths are tracked.

\paragraph{Mode crossing.} A value whose type contains no \code{List} anywhere---\code{Int},
\code{String}, a record of them, \code{Order}---can hold no cell and is never tracked. It
\emph{crosses} the ownership axis in the sense of Peters et al., for whom a type crosses an axis
``when all of its values have no interaction with that axis'' \citep[\S2.3]{peters2026mode}. We write
$\mathit{crosses}(\tau)$ when $\tau$ is a primitive, or a record, tuple or nominal type from which no
\code{List} is reachable through fields, aliases or constructor payloads. A type variable does not
cross: a value of type \code{a} may be a list, and it is tracked as a whole, at the single path
$\varepsilon$. Most values in a beni program are scalars and records of scalars, so crossing keeps
the analysis small at no cost in precision.

\subsection{Holders, and the two kinds}
\label{sec:model-holders}

A \emph{holder} of a cell $c$ at a program point is a pair $(\mathit{place}, \mathit{path})$
through which $c$ can be reached and read later. Table~\ref{tab:holders} lists the places. The
operand of an update is itself a holder; the check asks about all the others.

\begin{table}[t]
\centering\small
\begin{tabularx}{\linewidth}{@{}>{\raggedright\arraybackslash}p{0.34\linewidth}>{\raggedright\arraybackslash}p{0.2\linewidth}>{\raggedright\arraybackslash}X@{}}
\toprule
Place & Example & Kind \\
\midrule
a local or parameter & \code{xs}; \code{model.rows} as $(\mathit{model}, .\mathit{rows})$ & content \\
a field of a record, tuple or constructor held by a place & $(m, .\mathit{form}.\mathit{tags})$ & content \\
an element of a list held by a place & $(\mathit{rows}, [\ast].\mathit{tags})$ & content \\
a captured variable of a closure held by a place & $(f, \mathit{cap}\ \mathit{xs})$ & content, live while the closure is live \\
the environment of a suspended fiber, a command or subscription thunk, a mutable \code{Ref}, a \code{Queue} & \code{Cmd.perform (λ… → …xs…)} & content, \emph{escaped}: the analysis cannot see its future reads \\
a JavaScript sibling, or the result of \code{Js.from} & \code{Js.from xs} & content, escaped \\
a top-level value & \code{empty = []} & content, escaped for the program's life \\
a value the runtime retains and hands out again: a fiber's outcome (every joiner receives the same value), a bracketed resource, a queue's contents & \code{Task.join fib} & content; $\top$ on every path of the result \\
the module-level model between dispatches, and an event payload read from it or from a rendered row's item & \code{onClick=\{Keep model.rows\}} & content (\S\ref{sec:hard-tea}) \\
the runtime's per-row instance item & \code{insts[i].it} & identity only (A4) \\
a renderer slot & a derived value & leaf values only; a list-valued slot is recomputed before it is read (\S\ref{sec:hard-tea}) \\
\bottomrule
\end{tabularx}
\caption{The places that can hold a cell.}
\label{tab:holders}
\end{table}

\subsection{Unique at a site}
\label{sec:model-unique}

Let $u$ be an \emph{update site}: a call of an update primitive (a \emph{demand},
\S\ref{sec:model-demands}) whose operand expression is $e$. By beni's evaluation order the write
happens after every argument of the call has been evaluated. Let $C(e)$ be the set of cells $e$ may
evaluate to.

\begin{definition}[Dead; unique at $u$]
\label{def:unique}
A holder $h$ is \emph{dead} at a program point when no execution continuing from that point reads
through $h$. A \emph{read through $h$} is a dereference of the cell's elements or length, directly
or by passing $h$'s value to anything that may read it. A cell $c \in C(e)$ is \emph{unique at $u$}
when every holder of $c$---the operand's own binding included---is dead at $u$. The operand's
occurrence itself is not a later read: it was evaluated before $u$, and the result of the demand is
the new owner of the cell.
\end{definition}

\begin{definition}[The rule]
An update site $u$ is accepted if and only if every $c \in C(e)$ is unique at $u$, $\top \notin
C(e)$ (the checker knows which cells are involved), and no $c$ is \emph{escaped} (held by something
whose future reads the checker cannot see).
\end{definition}

Three consequences follow. First, \emph{reads before writes are never sharing}.
\code{List.set xs i (List.get xs j + 1)} reads \code{xs} while evaluating its third argument, before
the call; if nothing reads \code{xs} afterwards, the update is accepted. This is Wadler's design goal
that ``only update operations are sequentialised but lookups can occur in parallel''
\citep[\S3]{wadler1991use}, and the reservation phase of Rust's two-phase borrows, during which a
mutable borrow ``acts exactly like a shared borrow'' until it is activated \citep{rfc2025}.
Second, a \emph{later} read is sharing: in
\begin{Verbatim}
ys = List.set xs 0 1
List.get xs 0
\end{Verbatim}
the second line reads the old version, and the program is rejected with a message naming the
binding, the \code{set} and the \code{get}. Third, \emph{a dead alias is harmless}:
\code{saved = xs} followed by \code{List.set xs 0 1}, with \code{saved} never read again, is
accepted---Boyland's condition that aliases need only ``all be dead''
\citep[\S3.2]{boyland2001alias}.

This is not Clean's uniqueness. Clean's \code{*} attribute is a structural property---exactly one
reference exists---checked by counting occurrences \citep[Def.~8.14]{barendsen1996uniqueness},
refined by an observing-reference rule for guards and strict \code{let}s
\citep[\S9.4]{plasmeijer2011clean}. Liveness-based uniqueness is strictly more permissive and needs
no \code{let!} construct \citep{wadler1990linear}.

\subsection{Abstract values}
\label{sec:model-abstract}

For every expression the analysis computes an \emph{abstract value} $A(e)$: a finite map from the
list-typed paths under the value to sets of abstract cells,
\[
A(e) : \mathit{Path} \rightharpoonup \wp(\mathit{Cell})
\qquad
\begin{array}{r@{\;}c@{\;}l@{\quad}l}
\mathit{Cell} & ::= & \site(\iota) & \text{a list made at instruction } \iota \text{ (a literal, a fresh result)}\\
 & \mid & \pi_i.q & \text{the cell at path } q \text{ of parameter } i \text{, on entry}\\
 & \mid & \top & \text{a cell the analysis cannot name}
\end{array}
\]
We write $A(e)@q$ for the set at path $q$. Two conventions keep the domain honest. \emph{$\emptyset$
means provably no cell}---a path of a crossing type, or a fresh container known to be empty---and
never ``unknown''. \emph{$\top$ is prefix-closed}: a value whose path $q$ maps to $\top$ has $\top$
at every path under $q$. A decoder's \code{List (List Int)} whose outer cell is fresh but whose inner
cells are undescribed therefore has $[\ast] \mapsto \{\top\}$, and the result of a call with no
summary is $\top$ everywhere.

The analysis is \emph{directional}, in the style of Andersen's inclusion-based points-to analysis
\citep{andersen1994program}: a binding \code{y = x} sets $A(y) = A(x)$, a copy of the sets, never a
merged equivalence class as in unification-based analyses \citep{steensgaard1996points}. The
difference is the precision argument. Emre et al., studying the translation of C to safe Rust,
found that under a unification-based model ``having only four necessarily unsafe raw pointers is
enough to poison 85\,\% of the total raw pointers'' in tmux, and that a subset-based analysis raised
the share of well-contained pointers from 6.7\,\% to 88.9\,\% \citep[\S5,
Table~2]{emre2023aliasing}. Hindley--Milner unification is therefore the wrong solver for this fact;
beni's effect bits already live beside the unifier rather than in it, and ownership does the same.

\subsection{Transfer rules}
\label{sec:model-transfer}

Table~\ref{tab:transfer} gives the abstract value of each core form, in evaluation order, with
$\Sigma$ the environment of locals. On abstract values, $\cup$ is pointwise union,
$\mathit{shift}(A, s)$ prefixes every path with step $s$, and $\mathit{proj}(A, s)$ keeps the paths
under $s$ with the step removed.

\begin{table}[t]
\centering\small
\begin{tabularx}{\linewidth}{@{}>{\raggedright\arraybackslash}p{0.25\linewidth}>{\raggedright\arraybackslash}p{0.33\linewidth}>{\raggedright\arraybackslash}X@{}}
\toprule
Expression & $A(e)$ & Holders created; notes \\
\midrule
variable $x$ & $\Sigma(x)$ & none new; a \emph{use} for liveness \\
literal \code{[ x₁, …, xₙ ]} & $\varepsilon \mapsto \{\site(\iota)\} \cup \bigcup_i \mathit{shift}(A(x_i), [\ast])$ & a fresh cell; the elements' cells are reachable through $[\ast]$ \\
record, tuple, constructor & $\bigcup_f \mathit{shift}(A(x_f), .f)$ & the structure holds what its parts hold \\
record update \code{\{ r | f = y \}} & $A(r)$ with the paths under $.f$ replaced by $\mathit{shift}(A(y), .f)$ & the other fields are the very same values, so old and new records share every cell under every other field \\
field access \code{r.f} & $\mathit{proj}(A(r), .f)$ & an alias of the field's cell, not a copy \\
\code{case x of p → e} & per arm, pattern variables bound to projections of $A(x)$; the \code{rest} of \code{[ x, …rest ]} is bound to \emph{the same cells as $x$}; the \code{init} of \code{[ …init, last ]} to a fresh cell & arms are analysed separately; results are joined by union; liveness is per path \\
\code{x = e₁} then \code{e₂} & $\Sigma[x \mapsto A(e_1)]$ & $x$ holds everything in $A(e_1)$ \\
lambda \code{λy₁ … yₙ → e} & a function value with summary $S_\lambda$ (\S\ref{sec:model-summary}) & the closure holds each captured $z$ at $(\lambda, \mathit{cap}\ z)$ \\
call \code{f a₁ … aₙ} of a top-level $f$ & by the summary $S_f$ & a use of every argument; a \emph{demand} when some parameter is consumed \\
call \code{g a₁ … aₙ} through a value & by $g$'s ownership class (\S\ref{sec:inf-classes}) & as above \\
\code{foreign} or \code{Js} operation with no summary & $\varepsilon \mapsto \{\top\}$, prefix-closed & every argument's cells become escaped \\
top-level value $v$ & $\Sigma(v)$ & every cell escaped (\S\ref{sec:model-check}) \\
\bottomrule
\end{tabularx}
\caption{Transfer rules of the cell analysis.}
\label{tab:transfer}
\end{table}

\subsection{Summaries}
\label{sec:model-summary}

The compiler infers a summary for every function and publishes it in the module interface. For a
function $f$ with parameters $\pi_1 \ldots \pi_n$ and result $\rho$:
\[
\begin{array}{r@{\;}l@{\qquad}l}
S_f = \langle & \mathit{use} : (i, q) \mapsto \borrowed \mid \shared \mid \consumed & \text{for each list-typed parameter path } \pi_i.q\\
, & \mathit{res} : q \mapsto \wp(\{\pi_i.q'\} \cup \{\fresh\}) & \text{for each list-typed result path } \rho.q\\
, & \mathit{fun} : j \mapsto \langle \mathit{calls}, \mathit{escapes}, \mathit{supply}, \mathit{need} \rangle & \text{for each function-typed parameter } \pi_j\\
, & \mathit{reason} : (i, q) \mapsto \text{a derivation tag} & \text{why each non-borrowed use was inferred} \;\rangle
\end{array}
\]
with $\mathit{calls} \in \{0, {\le}1, \mathit{many}\}$ and $\mathit{escapes} \in \{\mathit{no},
\mathit{yes}\}$.

\begin{itemize}
\item \emph{Borrowed}: during the call $f$ may read the cell at $\pi_i.q$; it does not write it,
store it, return it, or capture it in anything that outlives the call. After the call the caller
still owns the cell. This is Aspinall, Hofmann and Kone\v{c}n\'y's third usage aspect, ``read and not
shared'' \citep[\S1]{aspinall2008usage}, and FP\textsuperscript{2}'s borrowed parameter, which
``cannot be used in a destructive match, passed as an owned parameter, or returned as a result''
\citep[\S1.3]{lorenzen2023fp2}.
\item \emph{Shared}: $f$ may make the cell reachable from its result (and $\mathit{res}$ says where),
or from something that outlives the call---a stored closure, a \code{Ref}, a fiber, JavaScript.
\item \emph{Consumed}: $f$ may write the cell in place, directly or by passing it to a consumed
position. The caller must prove the cell unique at the call. A consumed parameter may also appear
in $\mathit{res}$---every list writer returns its operand---so ownership is transferred, not lost.
\item $\mathit{res}$ records which arguments the result may alias, per path. It is the dependent
result qualifier of reachability types, where a function type names the argument its result
reaches, as in \code{id(x: T}$^{\blacklozenge}$\code{): T}$^{\{x\}}$ \citep[\S2.2]{wei2024polymorphic}.
\item $\mathit{fun}(j)$ records, for a function-typed parameter $g$, how often $f$ may call it (a
call inside a loop, a recursion or a callee that maps over a list is $\mathit{many}$), whether it
escapes (is stored, returned, or captured by something stored), what $f$ \emph{supplies} at each of
$g$'s parameters (an owned or a borrowed cell), and what $f$ \emph{needs} of $g$'s result (whether
it stores it, or relies on it being fresh). The lambda passed at a call is checked against these
facts (\S\ref{sec:inf-classes}). The ${\le}1$/$\mathit{many}$ distinction is OxCaml's
\code{once}/\code{many}: when a unique value is consumed in a closure, ``this closure can be
invoked at most once'' \citep{oxcaml2025uniqueness,lorenzen2024oxidizing}.
\end{itemize}

A lambda's summary $S_\lambda$ has the same shape over its own parameters, plus
$\mathit{cap} : (z, q) \mapsto \borrowed \mid \shared \mid \consumed$ for each captured variable
\emph{per path}, and a flag $\mathit{once}$, set exactly when some capture is consumed. Per-path
captures matter in TEA code: a command lambda that reads only \code{model.page} captures
$(\mathit{model}, .\mathit{page})$ and holds nothing of \code{model.todos}, which keeps the common
\code{( \{ model | todos = … \}, Cmd.perform (λsend → Http.get (url model.page) …) )} from sharing every
list in the model. \emph{A $\mathit{once}$ closure is an affine value}: its binding is dead after one
use on every path, and a second use---a second call, a second pass to any position, a store---is an
error. The side conditions for its consumed captures are discharged at the lambda's \emph{creation}
against the closure's own liveness. Thus
\begin{Verbatim}
g = λx → List.push acc x
r1 = Result.andThen a g
r2 = Result.andThen b g
\end{Verbatim}
is rejected at the second use of \code{g}, even though each \code{Result.andThen} calls its callback
at most once.

\paragraph{Conditional summaries.} Summaries may depend on the class of a function-typed parameter.
\code{List.foldl}'s accumulator is consumed only when its callback consumes its own accumulator
parameter: $\mathit{use}(\mathit{acc}) = \consumed$ iff $\mathit{need}_g(2) = \consumed$. A body whose
check depends on whether a closure handed to a parameter escapes, such as
\begin{Verbatim}
f xs h =
    g = λ_ → List.length xs
    _ = h g
    List.push xs 1
\end{Verbatim}
is accepted iff $\mathit{fun}_h(1).\mathit{escapes} = \mathit{no}$, and that condition travels in
$f$'s summary to be checked at each caller when $h$'s class is known. Formally, every $\mathit{use}$,
$\mathit{res}$ and side condition may be an expression over the class variables of the function's
function-typed parameters: rank-1 constraints, the same shape the effect summaries already have.

\subsection{Demands: declared, not inferred}
\label{sec:model-demands}

The list writers in beni's core library are written in beni over a low-level \code{Js} interface
(for example, \code{set} slices the array through \code{Js.call}). By the transfer rules, a value
passed to \code{Js} is escaped and a \code{Js} result is $\top$, so inference would conclude that every
writer ``escapes its operand and returns an unknown cell''. The writers therefore carry an
\emph{asserted} summary, a promise the body cannot show, just as a \code{foreign} declaration carries
an asserted effect rung. The core library and the platforms may assert a summary; every other
package may only \emph{narrow} an inferred one (\S\ref{sec:disc-options}). A false assertion is a
correctness bug in privileged code, not a user error. In bracket notation:
\begin{Verbatim}
set      : List a [consumed], Int, a → List a [= π₁]
update   : List a [consumed], Int, (a → a) [calls ≤1, supply owned if π₁ excl, need owned]
             → List a [= π₁]
push     : List a [consumed], a → List a [= π₁]
pop      : List a [consumed] → List a [= π₁]
swap     : List a [consumed], Int, Int → List a [= π₁]
insertAt : List a [consumed], Int, a → List a [= π₁]
removeAt : List a [consumed], Int → List a [= π₁]
cons     : a, List a [consumed] → List a [= π₂]
append   : List a [consumed], List a [borrowed] → List a [= π₁, fresh]
length   : List a [borrowed] → Int
copy     : List a [borrowed] → List a [fresh]          -- new
at       : List a [borrowed], Int → a [= π₁[*]]        -- core-private element read
\end{Verbatim}
The user-facing demand set is the \emph{nine writers} with a consumed parameter: \code{set},
\code{update}, \code{push}, \code{pop}, \code{swap}, \code{insertAt}, \code{removeAt}, \code{cons} and
\code{append}. (\code{update} is \code{set} with a callback.) Everything else in the list library is
ordinary beni over \code{at} and an array builder, and its summary is inferred: \code{map}, for
instance, comes out as
\code{List a [borrowed], (a → b) [calls many, supply borrowed, escapes no] → List b [fresh]}. Putting
the demand only on the update primitives follows a Clean recipe: library functions such as
\code{newArray} return an array whose attribute is a variable rather than ``unique'', because ``if
we are careful never to return a unique array from a function, we will always be able to share
arrays'' \citep[\S6]{devries2008simplified}. User code carries nothing. \code{append}'s result never names
$\pi_2$; this requires dropping one identity promise of the current library (\S\ref{sec:disc-options}).

\paragraph{A summary describes the emitted code.} beni's backend compiles a function that builds
its result in tail position modulo a list constructor, such as
\begin{Verbatim}
appendTo xs ys = case xs of
    [] → ys
    [ x, …rest ] → [ x, …appendTo rest ys ]
\end{Verbatim}
into a builder loop whose exit \emph{copies} \code{ys}, whereas the source's \code{cons} steps would
consume the result that holds \code{ys}. The checker reads such a function as the backend emits
it---the exit value is borrowed, the result fresh---because the rewrite is mandatory and part of the
language's cost contract. The same holds for loops whose list traversal the backend rewrites into
an offset into the base array: a slot so rewritten is a read of the base.

\subsection{Liveness, per path}
\label{sec:model-liveness}

$\live(h, p)$ holds for a holder $h = (x, q)$ at point $p$ when there is a use of $x$ at or after $p$,
on some path, that reads $q$ or a prefix of it. The uses, in evaluation order, are:
\begin{itemize}
\item a read of $x.q'$ with $q' \sqsubseteq q$ or $q \sqsubseteq q'$ (prefix order);
\item $x$ passed to a call whose summary reads the parameter at a path that meets $q$; a parameter
with no summary reads everything;
\item $x$ returned, stored in a structure that is live, or captured by a closure that is live---a
captured variable is live exactly while the closure is, and forever once the closure escapes;
\item $x$ used whole by anything the analysis does not see into (\code{Debug.toString}, a \code{Js}
operation): a read of every path. \code{==} and \code{<} are method calls: a derived \code{eq} or
\code{compare}, and the library's own list equality and comparison, are read-only by construction,
but a type's own \code{eq} is ordinary beni and is summarised like any function---one that consumes
is a demand.
\end{itemize}
A \code{case} makes liveness path-sensitive in the usual way: a holder live only in the arm not
taken is dead in the other. A self tail call rebinds the parameters, and the old values are dead
after the jump unless an argument still to be evaluated reads them.

Liveness is computed per (variable, path), so that \code{m.name} read after
\code{List.push m.rows x} is not a use of \code{m.rows}. Field-sensitive deadness is what lets a
record be threaded through helpers that each touch one field. Rust's borrow checker lacks this
across function boundaries, and in our prototype study its absence cost a copy of every row on each
append.

\subsection{The check}
\label{sec:model-check}

At a call $f\,\bar a$ whose summary has $\mathit{use}(i, q) = \consumed$, for each $c \in A(a_i)@q$:

\begin{description}[leftmargin=1.2em,font=\normalfont]
\item[\Known] $c \neq \top$, and $A(a_i)@q$ is not empty for lack of information.
\item[\Unescaped] $c$ is not escaped: nothing whose reads the checker cannot see holds it.
\item[\Dead] for every holder $h$ of $c$, $\neg\live(h, \text{after the call})$: nothing reads it later.
\item[\Disjoint] for every other argument path $(a_j, q')$ with $j \neq i$ or $q' \neq q$, and for every
capture of every function-valued argument, $c \notin A(a_j)@q'$: nothing reads it \emph{during} the call.
\end{description}

\Dead\ reads ``after the call'': every argument has been evaluated, so a read inside an argument is
before the write, and the operand's own binding is a holder like any other. \Disjoint\ is the
separation rule. The write happens \emph{inside} the callee, before the callee's own later reads,
and the callee reads its other parameters and calls its callbacks while the write is already visible.
In \code{List.foldl xs xs (λx acc → List.insertAt acc 0 x)} the walked input and the consumed
accumulator are the same array, so the fold would read the growing array through \code{xs} while
writing it through \code{acc}. In \code{g xs (λ\_ → List.length xs)} with
\code{g xs k = k (List.push xs 1)}, the callback would read the pushed array through its capture.
\Disjoint\ corresponds to Rust's rule of one mutable or any number of shared borrows at a call
\citep{rfc2094}, to Futhark's requirement that the argument for a consumed parameter not alias any
other argument \citep{henriksen2017thesis}, and to Aspinall, Hofmann and Kone\v{c}n\'y's reading of the
comma in a typing context as a tensor \citep[\S3]{aspinall2008usage}. It is checked on abstract
cells, so two arguments that \emph{may} be one cell fail it.

When $a_i$ is itself a parameter path $\pi_j.q'$ of the enclosing function, that function's
$\mathit{use}(j, q')$ becomes $\consumed$, and \Unescaped, \Dead\ and \Disjoint\ are discharged at
\emph{its} callers in turn. A failure of \Known\ or \Unescaped\ is reported as a \emph{limit} of the
checker; a failure of \Dead\ or \Disjoint\ as \emph{sharing}, with the holder named
(\S\ref{sec:errors}).

\paragraph{Holders that live for the whole program.} A top-level value, including one with a
\code{where} clause that is memoised per evidence, is a holder that is never dead: its cells are
escaped from the start, so \code{[ x, …defaults ]} onto a top-level list is rejected, with
\code{List.copy} as the fix. The page's model variable, between dispatches, is the other such holder
(\S\ref{sec:hard-tea}).

\subsection{Deep uniqueness: exclusive contents}
\label{sec:model-excl}

A list's elements, a dictionary's values and a record's fields can hold cells, and two positions of
one container can hold the \emph{same} cell: \code{[ r, r ]}, \code{List.repeat r 3}, or a dictionary
into which one list was inserted under two keys. An edit inside an element---
\code{\{ r | tags = List.push r.tags t \}} on the element at position 0---would then be seen at
position 1. The cell-set domain cannot tell positions apart: every element of a list built in a loop
comes from one allocation site, so $A(\mathit{rows})@[\ast].\mathit{tags}$ is one abstract cell for
all of them, and \Dead\ finds the holder $(\mathit{rows}, [\ast].\mathit{tags})$ live whenever the
container outlives the edit. That is the right answer when positions may share and a false error
when they cannot. It is the case Clean handles with \emph{deep} uniqueness, a unique list of unique
elements \citep[\S9.2]{plasmeijer2011clean}; OxCaml likewise treats uniqueness as ``a deep property
by default'' \citep[\S2.1]{lorenzen2024oxidizing}.

We add a flag $\excl$ to abstract container paths---$(x, q[\ast])$ for a list, $(x,
q.\mathit{contents})$ for a dictionary. A container path is \emph{exclusive} when every cell stored
under it was unique at the store and stored nowhere else, and no read since has produced a holder of
a contained cell that is still live. The flag is
\begin{itemize}
\item \emph{set} by a construction whose every stored cell is unique at the store: a literal of
distinct fresh cells, or a fresh result whose elements are declared fresh (\code{map}'s result when
the callback's result is fresh; a decoder's output, by its declaration);
\item \emph{kept} by a \emph{destructive read}: a container function that consumes the container,
hands one element to a callback and stores the callback's result back at that position, reading no
other element in between. \code{List.update}'s body is \code{set xs i (func (at xs i))}, and
\code{set} writes slot $i$. The element is unique inside the callback because the container is
consumed at the call (\Dead), no argument or capture of the callback holds it (\Disjoint), and so the
slot is the element's only other holder and nothing reads the slot before it is overwritten. A
dictionary's \code{update} has the same shape at a \emph{key}, which the path domain cannot name; its
callback's $\mathit{supply} = \mathrm{owned}$ is therefore an asserted row of the core library;
\item \emph{cleared} by a store of a cell that is not unique at the store (the second \code{r} of
\code{[ r, r ]}, \code{List.repeat}, an element taken from another live container), and by a
non-destructive read whose result is live while the container is (a dictionary lookup whose
dictionary is used afterwards, \code{List.get}, the \code{x} of a pattern \code{[ x, …rest ]} while
the scrutinee lives).
\end{itemize}
With the flag, \code{List.update rows i (λr → \{ r | tags = List.push r.tags t \})} is accepted when
\code{rows} is unique and exclusive, and rejected as a limit---``the elements of \code{rows} may
share''---when they came from somewhere the checker cannot see. The flag travels in summaries beside
$\mathit{res}$ (a result path may be $\fresh$ and $\excl$) and in interfaces. Elements of a crossing
type need none of this.
